\section{Our construction of layouts}\label{sec:layouts}

We build on top of Definition~\ref{def:layout}. First we define the spaces.

\begin{definition}[Spaces]\label{def:spaces}
Each side of a layout carries a space from $\{\logical, \physical, \tv\}$,
written $A \to B$. A layout may be annotated $A \to B$ for
$A, B \in \{\logical, \tv\}$, or $\logical \to \physical$, and no other pair is
constructible. Composition takes $A \to \logical$ and $\logical \to C$ to
$A \to C$, so $\tv \to \physical$ exists only as a composite.
\end{definition}

\begin{definition}[Replication]\label{def:linear}
Extend Definition~\ref{def:layout} by allowing the stride against a positive
integer $n$ to be the mark $\bot$ in place of an integer, with
\[
(n{:}\bot)(i) = 0 \quad\text{for every } i \in [0,n).
\]
\end{definition}

A zero stride is rejected rather than read as replication. This is because
replication is legal for reads and not for writes, since simultaneous loads of
one address are a broadcast while simultaneous stores are a race. It is better
not to use $0$, so that the two can never be mixed.

\begin{definition}[Swizzled layout]\label{def:swizzle}
A \emph{swizzle} is a bijection $\physical \to \physical$ of the form
\[
\sigma(x) \;=\; x \oplus \big(((x \gg \mathit{src}) \wedge (2^{b}-1)) \ll
\mathit{dst}\big),
\]
its two fields $[\mathit{src},\mathit{src}{+}b)$ and
$[\mathit{dst},\mathit{dst}{+}b)$ disjoint. A layout $A \to \physical$
followed by a swizzle is a \emph{swizzled layout}, again $A \to \physical$.
\end{definition}

A swizzle runs from $\physical$ to $\physical$, so it can only ever be the last
step, and being a bijection it sends distinct addresses to distinct addresses.
It therefore leaves the conditions below untouched, and matters only in the
arithmetic that is finally emitted.

\begin{definition}[Composition]\label{def:compose}
Let $f : A \to \logical$ and let $g : \logical \to C$ have domain shape $S_g$.
The \emph{composite} $g \circ f : A \to C$ is defined when $f$ is a dense
bijection onto $[0, \size S_g)$, and is the map
\[
g \circ f \;=\; c \mapsto g\big(\mathrm{unflatten}_{S_g}(f(c))\big),
\]
whose domain shape is that of $f$.
\end{definition}

A chain is therefore one $\tv \to \logical$ map, then any number of
$\logical \to \logical$ maps, then one $\logical \to \physical$ map.

\begin{definition}[Declared replicas]\label{def:replica}
For a layout $L$ with domain shape $S$, define a relation $\sim_L$ on
$\mathrm{Coord}(S)$:
\begin{itemize}
\item if $L$ is a layout, $c \sim_L c'$ when $c$ and $c'$ agree at every
mode whose stride is an integer, that is, they differ only at modes
with stride $\bot$;
\item if $L = g \circ f$ with $g$ of domain shape $S_g$, then $c \sim_L c'$ when
$\mathrm{unflatten}_{S_g}(f(c)) \sim_g \mathrm{unflatten}_{S_g}(f(c'))$.
\end{itemize}
\end{definition}

\begin{definition}[Validity]\label{def:validity}\label{sec:validity}
Let $L : A \to B$ be a layout with domain shape $S$. Then $L$ is \emph{valid}
when
\[
\begin{cases}
L \text{ is a dense bijection}
  & \text{if } B \in \{\logical, \tv\},\\[2pt]
L(c) = L(c') \implies c \sim_L c'
  & \text{if } B = \physical.
\end{cases}
\]
It is \emph{write-valid} when $B = \physical$ and in addition
\[
L(c) = L(c') \implies c = c'.
\]
Only a write-valid layout may be the target of a store.
\end{definition}

Validity admits overlap only as the replication of whole modes: two
coordinates may share an address only if they differ at modes of stride
$\bot$ alone. A map whose images overlap in part, such as the sliding windows
of a convolution, where neighbouring windows share some elements and not
others, is therefore not valid. Such read-only overlaps are outside the scope
of this paper.

When a layout's output is a logical or thread-value index, being valid means
being a dense bijection. This need not be checked by enumerating coordinates:
whether a layout is a dense bijection depends only on its sizes and strides.

\begin{lemma}[Dense bijections]\label{lem:dense}
Let $L$ be a flat layout with no stride $\bot$, its modes carrying sizes
$n_1,\dots,n_k$ and strides $s_1,\dots,s_k$ with every $n_i \ge 2$, and put
$N = \prod_i n_i$. Then $L$ is a dense bijection onto $[0,N)$ if and only if,
after sorting the modes by stride, $s_{(1)} = 1$ and
$s_{(j+1)} = s_{(j)}\,n_{(j)}$ for every $j$: the sorted strides are the
running products of the sorted sizes.
\end{lemma}

The two restrictions cost nothing. A mode of size $1$ only ever takes the
index $0$, so it contributes nothing and is dropped beforehand; a nested layout
is read through the flat layout of its leaves.

\begin{proof}
($\Leftarrow$) With running-product strides the map is the flattening of
Definition~\ref{def:layout} for the shape $(n_{(k)},\dots,n_{(1)})$, a
bijection onto $[0,N)$.

($\Rightarrow$) No stride is negative: index $1$ at a mode of negative stride,
and $0$ elsewhere, would put a negative value in the image. Sort the modes by
stride and write $N_j = \prod_{i \le j} n_{(i)}$
(so $N_0 = 1$). \emph{Claim:} $s_{(j)} \le N_{j-1}$ for every $j$. Every image
value below $s_{(j)}$ comes from a coordinate supported on modes
$1,\dots,j{-}1$, since a nonzero index at any later mode contributes at
least $s_{(j)}$ and no contribution is negative; those modes produce at
most $N_{j-1}$ distinct values, so surjectivity onto
$[0,\min(s_{(j)},N))$ forces $\min(s_{(j)},N) \le N_{j-1}$, and as
$N_{j-1} < N$ this is $s_{(j)} \le N_{j-1}$. Consequently
\[
\textstyle\sum_j (n_{(j)}-1)\,s_{(j)} \;\le\; \sum_j (n_{(j)}-1)\,N_{j-1}
\;=\; N - 1,
\]
the last equality by telescoping. But the left side is the maximum of the
image, and surjectivity onto $[0,N)$ requires the maximum to be exactly
$N-1$. The two sums therefore agree, and since each term on the left is at most
the corresponding term on the right, they agree termwise; as
$n_{(j)} - 1 \ge 1$ this gives $s_{(j)} = N_{j-1}$ for every $j$, which is the
running-product condition.
\end{proof}
